% Chapter 3, Figure 9: a fluid element in Cartesian coordinates, after Schlichting (scan: figures/ch3/fig09.png,
% PDF 325). Both parts are drawn in the house orthographic view (elevation 32 deg, azimuth 45 deg, tamr view=1.75),
% TikZ x = the flow direction x (up to the right), TikZ z = y (up), TikZ y = -z (the book's axes right-handed;
% dz comes towards the viewer, down to the right), so that the velocity arrows of the profile and the shear
% stresses on the element, both along x, are parallel, as the 1973 art draws them (its oblique view keeps x
% horizontal; an orthographic view with three faces of the element seen cannot).
% Left: the velocity profile u(y) in a plane parallel to the element's x-y plane (its base line u = 0 along y,
% its arrows along x), a qualitative sketch matched to the 1973 art (fig09.py: the traced curve fitted by a
% cubic, fig09.csv; the 18 velocity arrows from the base line to the curve, fig09-arrows.csv), its height
% \Hy = 2.3 times the element's dy (the art: 2.2). The profile in the Chapter 3 profile family's styles (Figs 6,
% 9, 16, 17, 18, 25): the profile u(y) as a data curve (s1, 1pt), its velocity arrows s1, lighter, with small
% heads; the base line an outline.
% Right: the fluid element, a rectangular box dx by dy by dz, drawn with dx = 1.2, dy = 1, dz = 0.75 (the 1973
% art draws a cube in an oblique view: a cube in this view would put the face diagonals and both face centres
% on the front vertical edge). The diagonals of the upper and lower y-faces (the lower ones, and the three
% edges at the far bottom corner, hidden) locate the face centres; from them the shear stresses on the two
% y-faces, tau + (partial tau/partial y) dy along +x on the upper face, tau along -x on the lower face
% (eq. (13)). The lower arrow is hidden inside the element until it leaves it under the near x-face (the 1973
% art draws it through the front face). Lettering as printed.
\documentclass[11pt]{standalone}
\usepackage{tamrfig}
\tikzset{
  profile/.style={draw=s1, line width=1pt},
  profile arrow/.style={draw=s1, line width=0.5pt, -{Stealth[length=3.4pt, width=2.4pt]}},
}
\begin{document}
\begin{tikzpicture}
  \def\vs{1.75}  % the view's scale (tamr view=\vs)
  \def\Hy{2.3}   % height of the profile, in the element's units (dy = 1)
  % ---- the velocity profile, in the house view: u along x, y along z (in units of its height) ----
  \begin{axis}[hide axis, disabledatascaling, anchor=origin, at={(0,0)},
      x={({\vs*\Hy*1.15cm},{\vs*\Hy*0.61cm})}, y={(0cm,{\vs*\Hy*1.38cm})},
      xmin=0, xmax=0.4, ymin=0, ymax=1, clip=false]
    \addplot[profile arrow, quiver={u=\thisrow{u}, v=0}]
      table[x expr=0, y=y, col sep=comma] {figures/v2/ch3/fig09-arrows.csv};
    \addplot[profile, no markers] table[x=u, y=y, col sep=comma] {figures/v2/ch3/fig09.csv};
    \draw[outline] (axis cs:0,0) -- (axis cs:0,1);
    \node[anchor=west] at (axis cs:0.37,0.5) {$u(y)$};
  \end{axis}
  % ---- the fluid element ----
  \begin{scope}[shift={(3.95,1.4)}, tamr view=\vs]
    \def\dx{1.2}\def\dy{1}\def\dz{0.75}
    \def\Lt{1.4}\def\Lb{1.3}   % the shear stress arrows' lengths from the face centres
    \coordinate (ct) at ({\dx/2},{-\dz/2},\dy);   % centre of the upper y-face
    \coordinate (cb) at ({\dx/2},{-\dz/2},0);     % centre of the lower y-face
    % the three seen faces (near x, near z, upper y)
    \blk{0}{\dx}{-\dz}{0}{0}{\dy}
    % hidden: the three edges at the far bottom corner (dx, 0, 0), the lower face's diagonals, and tau
    % inside the element, from the centre to the near x-face
    \draw[hidden] (0,0,0) -- (\dx,0,0) -- (\dx,-\dz,0) (\dx,0,0) -- (\dx,0,\dy);
    \draw[hidden] (0,0,0) -- (\dx,-\dz,0) (\dx,0,0) -- (0,-\dz,0);
    \draw[hidden] (cb) -- (0,{-\dz/2},0);
    % the upper face's diagonals
    \draw[thin line] (0,0,\dy) -- (\dx,-\dz,\dy) (\dx,0,\dy) -- (0,-\dz,\dy);
    % the shear stresses
    \draw[vec] (0,{-\dz/2},0) -- ({\dx/2-\Lb},{-\dz/2},0)
      node[anchor=south east, inner sep=1pt, xshift=1pt] {$\tau$};
    \draw[vec] (ct) -- ({\dx/2+\Lt},{-\dz/2},\dy)
      node[anchor=west, inner sep=2pt] {$\tau + \dfrac{\partial\tau}{\partial y}\,dy$};
    \node[point] at (ct) {};
    \node[point] at (cb) {};
    % the edges' lengths
    \node[anchor=north west, inner sep=2pt] at ({\dx/2},-\dz,0) {$dx$};
    \node[anchor=west, inner sep=3pt] at (\dx,-\dz,{\dy/2}) {$dy$};
    \node[anchor=north east, inner sep=2pt] at (0,{-0.8*\dz},0) {$dz$};
  \end{scope}
\end{tikzpicture}
\end{document}
